\section{Discussion}\label{sec:discussion}

\emph{When the layer acts.} The direction of the effect follows two
properties of the pipeline rather than the dataset: whether the tracker has a
continuation stage, and whether the detector's score distribution matches the
tracker's thresholds. On MOT17 with well-matched detections, the two-stage
trackers BoostTrack, ByteTrack at floor 0.1, and Hybrid-SORT produced output
identical to their hosts, while the single-stage OC-SORT and PD-SORT gained
HOTA on every sequence. When the stream fits the tracker, the layer stays out
of the way, and when the tracker lacks a stage, the layer supplies it.

\emph{Self-limiting by design.} In the clean regime the layer keeps the
tracker's thresholds and scores, so a tracker with a low-score stage and
well-matched detections receives its native input. This is the main
difference from the prior design \aob{}, which lowered the HOTA of the same strong
trackers by 4.15 to 5.89 points (Table~\ref{tab:prior}). The layer therefore does not
need to know in advance whether a tracker needs help; it reads the answer from
the score stream.

\emph{A continuation stage for single-stage trackers.} Single-stage trackers
drop every detection below one threshold, so a track can be lost when its
object scores low for a few frames. The continuation rescue considers only
candidates that overlap a track of the previous frame that no usable candidate
covers, so it targets exactly these gaps. OC-SORT and the external PD-SORT
gained HOTA on 7 of 7 MOT17 sequences through this stage.

\emph{Detector replacement.} When a detector is replaced or its confidence
scale shifts, the fixed thresholds of a tracker stop matching the scores. The
layer re-anchors them from the score stream without labels. Under halved
scores it restored two trackers from 0 to about 66 HOTA, and with a different
detector it raised the MOTA of two two-stage trackers by 6.0 and 7.9 points.
This is the situation the layer was built for: a detector that is swapped,
retrained, or quantized while the tracker stays fixed.

\emph{Real-time deployment.} The layer needs no training, no GPU, and no
access to the detector's internals, and it needs only four contract values
from the tracker's configuration. At run time it reads the candidate list, one
small grayscale image for the motion cue, and the tracks of the previous
frame. On a four-core CPU it takes 2.95 ms per frame, adds 4.5\% to the end-to-end time, and makes the tracker
itself faster. \scl{} therefore adds very little latency and is
suitable for real-time use. New trackers are added through their published
contract values and an adapter, without changes to the policy.
